\chapter{Realizacja systemu}
\label{chap:implementation}
\statusnote{Rozdział przeznaczony do uzupełniania po wykonaniu i zweryfikowaniu kolejnych komponentów. Obecnie repozytorium zawiera szkielet pracy i narzędzia do jej budowania.}
\section{Środowisko i organizacja repozytorium}
\writingnote{Opisz wersje narzędzi, strukturę kodu, Git, automatyczne testy i kompilację pracy. Podawaj sprawdzone instrukcje uruchomienia.}
\section{Symulator wielu węzłów}
\writingnote{Po implementacji opisz generator normalnej aktywności, niezależne tożsamości, konfigurację możliwości, seed i czas logiczny.}
\section{Kontrolowane scenariusze anomalii}
\writingnote{Udokumentuj parametry, czas trwania, limity i etykietowanie scenariuszy. Etykiety nie mogą stawać się cechami wejściowymi detektora.}
\section{Kolektor, baza danych i API}
\writingnote{Przedstaw przepływ wiadomości, walidację i najważniejsze testy integracyjne. Pokaż obsługę błędów i odłączenia brokera.}
\section{Firmware i integracja fizycznego prototypu}
\writingnote{Uzupełnij po uruchomieniu ESP32. Udokumentuj konfigurację GPIO, poziomy napięć, zasilanie napędów i rzeczywisty zakres sprzężenia zwrotnego.}
\section{Integracja uczenia maszynowego}
\writingnote{Po implementacji opisz przygotowanie cech, aktywację modelu, wynik detekcji, wersje i mechanizm wycofania modelu.}
